\centerline{\bf \Huge Комби покрытия}

\begin{flushright}
        \scriptsize{but why?}
\end{flushright}

\problem{1}
По кругу стоят 200 студентов из 10 групп, в каждой из которых 20 студентов. Докажите, что в каждой группе можно выбрать старосту таким образом, чтобы никакие два старосты не стояли бы рядом.

\problem{2}
Каждая клетка доски $100 \times 100$ окрашена либо в черный, либо в белый цвет, причем все клетки, примыкающие к границе доски - черные. Оказалось, что нигде на доске нет одноцветного клетчатого квадрата $2 \times 2$. Докажите, что на доске найдется клетчатый квадрат $2 \times 2$, клетки которого раскрашены в шахматном порядке.

\problem{3}
Из бумажного клетчатого квадрата $100 \times 100$ вырезали по границам клеток 1950 доминошек. Докажите, что из оставшейся части можно вырезать четырехклеточную фигуру «Т»-тетрамино (возможно, повернутую). (Если такая фигурка уже есть среди оставшихся частей, считается, что ее получилось вырезать.)

\problem{4}
По каждому из 100 видов работ в фирме имеется ровно 8 специалистов. Каждому сотруднику можно дать выходной в субботу или воскресенье. Докажите, что это можно сделать так, чтобы и в субботу, и в воскресенье для каждого вида работ присутствовал бы свой специалист. (Сотрудник может быть специалистом по нескольким видам работ; распределение специалистов по видам работ известно тому, кто назначает выходные.)

\problem{5}
Существует ли 2022-значное число, перестановкой цифр которого можно получить 2022 разных 2022-значных полных квадратов?

\problem{6}
Пусть $A$ есть 101-элементное подмножество множества $S=\left\{1,2, \ldots, 10^6\right\}$. Докажите, что для некоторых натуральных $t_1, \ldots, t_{100} \in S$ следующие множества $A_j=\left\{x+t_j\right.$ : $x \in A\}, j=1, \ldots, 100$, попарно не пересекаются.

\problem{7}
21 мальчик и 21 девочка участвуют в математической олимпиаде. Оказалось, что каждый участник решил не более 6 задач. Кроме того, для каждой пары, состоящей из мальчика и девочки, есть задача, которую они оба решили. Докажите, что есть задача, которую решили не менее трех мальчиков и не менее трех девочек.

\problem{8}
Можно ли выбрать 2022 различных натуральных чисел, не превосходящих 100000 , таких, что никакое из них не было бы равно среднему арифметическому двух других?